% META: {"id": "1SPE-GEOREP-COURS-01", "chapitre": "1SPE-GEOMETRIE-REPEREE", "type_objet": "cours", "sous_type": "diagnostic", "titre": "Diagnostic d'entree", "temps_gabarit": 2, "duree_min": 20, "prerequis_testes": ["R1", "R2", "R3", "R4", "R5"], "mode_creation": "ex_nihilo", "sources_inspiration": [], "status": "generated"}
% BEGIN-VERIFY
% from sympy import symbols, sqrt, Rational, solve, Symbol
% x, y = symbols('x y')
% # Q1 : distance entre A(1,3) et B(4,7)
% d = sqrt((4-1)**2 + (7-3)**2)
% assert d == 5
% # Q2 : milieu de A(2,6) et B(8,4)
% mx, my = Rational(2+8, 2), Rational(6+4, 2)
% assert (mx, my) == (5, 5)
% # Q3 : vecteur AB avec A(1,2) B(4,6)
% assert (4-1, 6-2) == (3, 4)
% # Q4 : colinearite u(2,3) v(4,6)
% assert 2*6 - 3*4 == 0
% # Q5 : produit scalaire u(3,-1) v(2,5)
% assert 3*2 + (-1)*5 == 1
% # Q6 : orthogonalite u(4,3) v(-3,4) : 4*(-3)+3*4=0
% assert 4*(-3) + 3*4 == 0
% # Q7 : discriminant de 2x^2 - 3x + 1
% a, b, c = 2, -3, 1
% delta = b**2 - 4*a*c
% assert delta == 1
% # Q8 : systeme 2x + y = 5, x - y = 1
% sol = solve([2*x + y - 5, x - y - 1], [x, y])
% assert sol == {x: 2, y: 1}
% END-VERIFY

\section*{Temps 2 — Diagnostic d'entrée \hfill \normalfont\small(15–20 min, sans calculatrice)}

Réponds à chaque question sur ta copie. Ne cherche pas à aller vite : l'objectif est de repérer ce que tu maîtrises, pas de tout réussir.

\bigskip

\subsection*{R1 — Repérage dans le plan}

\begin{enumerate}

  \item Calcule la distance entre les points $A(1 ; 3)$ et $B(4 ; 7)$.

  \item Détermine les coordonnées du milieu $M$ du segment $[AB]$ avec $A(2 ; 6)$ et $B(8 ; 4)$.

\end{enumerate}

\subsection*{R2 — Vecteurs et coordonnées}

\begin{enumerate}[resume]

  \item Détermine les coordonnées du vecteur $\vec{AB}$ avec $A(1 ; 2)$ et $B(4 ; 6)$.

  \item Les vecteurs $\vec{u}(2 ; 3)$ et $\vec{v}(4 ; 6)$ sont-ils colinéaires ? Justifie.

\end{enumerate}

\subsection*{R3 — Produit scalaire}

\begin{enumerate}[resume]

  \item Calcule le produit scalaire $\vec{u} \cdot \vec{v}$ avec $\vec{u}(3 ; -1)$ et $\vec{v}(2 ; 5)$.

  \item Les vecteurs $\vec{u}(4 ; 3)$ et $\vec{v}(-3 ; 4)$ sont-ils orthogonaux ? Justifie.

\end{enumerate}

\subsection*{R4 — Équations du second degré}

\begin{enumerate}[resume]

  \item Calcule le discriminant de l'équation $2x^2 - 3x + 1 = 0$ et déduis-en le nombre de solutions.

\end{enumerate}

\subsection*{R5 — Systèmes d'équations}

\begin{enumerate}[resume]

  \item Résous le système $\begin{cases} 2x + y = 5 \\ x - y = 1 \end{cases}$.

\end{enumerate}

\bigskip

\subsection*{Bilan}

\begin{center}
\begin{tabular}{|c|c|c|}
\hline
\textbf{Prérequis} & \textbf{Questions} & \textbf{Fiche si besoin} \\
\hline
R1 — Repérage & Q1, Q2 & Fiche R1 \\
\hline
R2 — Vecteurs & Q3, Q4 & Fiche R2 \\
\hline
R3 — Produit scalaire & Q5, Q6 & Fiche R3 \\
\hline
R4 — Second degré & Q7 & Fiche R4 \\
\hline
R5 — Systèmes & Q8 & Fiche R5 \\
\hline
\end{tabular}
\end{center}
